
\section{功~动能~动能定理}

\subsection{功的相关概念}

\subsubsection{机械功}

\textbf{力和物体在力的方向上移动距离的乘积，叫做力对物体做的功。}

	$$\mathrm{d}A=\mathbf{F}\cdotp\mathrm{d}\mathbf{r}=|\mathbf{F}| |\mathrm{d}\mathbf{r}|\cos\theta=F\mathrm{d}s\cos\theta$$
	
	$$ A=\int\mathrm{d}A=\int_{a}^{b}\mathbf{F}\cdotp\mathrm{d}\mathbf{r}=\int_{a}^{b} F\cos\theta\mathrm{d}s$$
	
关于概念，辨析以下几点：

1.力在时间上的累积作用是动量，力在空间上的累积作用是功。

2.功是标量，无方向，正负代表做功的正负，单位为J或$N\cdotp m$。

3.$\mathbf{F}$是恒力，对于变力，在一段极小的位移中仍可视作恒力。

4.由于力是矢量，可以在$x,y,z$轴进行分解，因此功也可以在$x,y,z$轴分解，总功等于三者的代数和。

\subsubsection{功率}

\textbf{力在单位时间内做的功，叫做功率。}
$$P=\dfrac{\mathrm{d}A}{\mathrm{d}t}=\dfrac{\mathbf{F}\cdotp\mathrm{d}\mathbf{r}}{\mathrm{d}t}=\mathbf{F}\cdotp\mathbf{v}$$

1.功率表示力的做功快慢程度。

2.功率的单位为W或$J/s$。

\subsection{能量}

\textbf{能量是物体状态的单值函数，是各种运动状态的一般量度。}

机械运动状态由位置和速率描述，因此位置和速率是机械运动的状态参量，位置决定势能，速率决定动能，机械能由动能和势能组成。

\subsection{动能定理}

$$\mathrm{d}A=\mathbf{F}\cdotp\mathrm{d}\mathbf{r}=F\cos\theta|\mathrm{d}r| $$
由牛顿第二定律，$$F\cos\theta=m\mathbf{a}_t=m\dfrac{\mathrm{d}\mathbf{v}}{\mathrm{d}t} $$
得到$$\mathrm{d}A=F\cos\theta|\mathrm{d}r|=m\dfrac{\mathrm{d}\mathbf{v}}{\mathrm{d}t}\mathbf{v}\mathrm{d}t=m\mathbf{v}\mathrm{d}\mathbf{v} $$
因此，$$A=\int_{\mathbf{v}_a}^{\mathbf{v}_b} m\mathbf{v}\mathrm{d}\mathbf{v}=\frac{1}{2}m{\mathbf{v}_b}^2-\frac{1}{2}m{\mathbf{v}_a}^2$$
$$A=E_{kb}-E_{ka}$$

1.动能定理连接起了功与能之间的关系，将过程量与状态量相联系，\textbf{合外力对质点做的功等于质点动量的增量}。

2.由于运用了牛顿第二定律，因此动能定理只能在惯性系下成立，但在不同的惯性参考系中，动能定理形式上保持不变，即形式上与参考系的选择无关。